\chapter{Analyse du métier et des données}

\section{Le cadre de Geofoncier}
\subsection{Rôle et fonctionnement du portail}
% Décrire Geofoncier, à quoi ça sert, qui l’utilise, quels types de versements.

\subsection{Spécifications techniques des données attendues}
% Formats attendus, contraintes qualité, règles de nommage/tri, métadonnées, etc.

\section{Typologie des archives à traiter}
\subsection{Registres de communes et documents associés}
% Format, état, variations, structure des pages, présence de tableaux/colonnes.

\subsection{Plans cadastraux et pièces annexes (si pertinent)}
% À garder seulement si ton projet les traite réellement.

\subsection{Spécificités de l'écriture et difficultés}
% Manuscrit, bruit, encre, rotation, pliures, contraste, etc.

